\documentclass[UTF8]{ctexart}
\usepackage{amsmath,amssymb}
\usepackage[a4paper,margin=2.5cm]{geometry}
\title{数学分析 0917}
\date{}

\begin{document}
\maketitle

\subsection*{0. Recall}

有理数集 $\mathbb Q$ 是稠密的：
\[
\forall a,b\in\mathbb Q,\quad
a<b\ \Longrightarrow\
\exists c\in\mathbb Q,\quad a<c<b.
\]
进一步，也可以证明任意两个不同的实数之间都有有理数。

无限集合 $S$ 称为\textbf{可数的}，是指存在双射
$\phi:\mathbb Z^+\to S$；即可以将其元素依次编号为
$s_1,s_2,s_3,\ldots$。

\subsection*{1. 有理数的可数性}

\subsubsection*{整数可数}

将整数写成两列：
\[
\begin{array}{cccc}
1&2&3&\cdots\\
0&-1&-2&\cdots
\end{array}
\]
交替从两列取数，可得
\[
1,0,2,-1,3,-2,\ldots
\]
因此 $\mathbb Z$ 可数。

\subsubsection*{有理数可数}

先将正有理数按分子、分母排成表：
\[
\begin{array}{cccc}
\frac11&\frac21&\frac31&\cdots\\[2pt]
\frac12&\frac22&\frac32&\cdots\\[2pt]
\frac13&\frac23&\frac33&\cdots\\
\vdots&\vdots&\vdots&
\end{array}
\]
沿副对角线逐条枚举，并跳过与已列数相同的分数，
便可列出所有正有理数。再交替列出正有理数、负有理数，
并加入 $0$，可知 $\mathbb Q$ 可数。

\subsubsection*{实数不可数}

反设 $(0,1)$ 中的所有实数可以列为
\[
\begin{aligned}
c_1&=0.a_{11}a_{12}a_{13}\cdots,\\
c_2&=0.a_{21}a_{22}a_{23}\cdots,\\
c_3&=0.a_{31}a_{32}a_{33}\cdots,\\
&\ \vdots
\end{aligned}
\]
统一选取不以无限个 $9$ 结尾的小数表示。
构造 $d=0.d_1d_2d_3\cdots$，其中
\[
d_i=
\begin{cases}
1,&a_{ii}\ne1,\\
2,&a_{ii}=1.
\end{cases}
\]
则 $d$ 的第 $i$ 位与 $c_i$ 的第 $i$ 位不同，
因此 $d$ 不在这张表中，矛盾。
所以 $\mathbb R$ 不可数。

由于 $\mathbb Q$ 可数，而 $\mathbb R$ 不可数，
无理数集也不可数。

\subsection*{2. 有理数没有连续性}

\subsubsection*{回顾：$\mathbb Q$ 的稠密性}

有理数集 $\mathbb Q$ 具有稠密性，即
\[
\forall a,b\in\mathbb Q,\ a<b,\ \exists c\in\mathbb Q,\quad a<c<b .
\]
（练习：证明 $\forall a,b\in\mathbb R,\ a<b$，
$\exists c\in\mathbb Q$，使得 $a<c<b$。）

\subsubsection*{直观例子}

取 $\pi$ 的逐位截断值：
\[
a_0=3,\quad a_1=3.1,\quad
a_2=3.14,\quad\ldots
\]
显然
\[
a_0<a_1<a_2<\cdots<a_n<\cdots,
\]
这些数均为有理数，且
\[
\lim_{n\to\infty}a_n=\pi\notin\mathbb Q.
\]
这提示我们：有理数的稠密性并不等于连续性。
（类似地，$\sqrt2=1.414\ldots$ 也总可以写成某个有理数
点列的极限，而它本身不属于 $\mathbb Q$。）

\subsubsection*{分割的定义}

设 $S$ 为有序集。若非空集合 $A,B$ 满足
\[
A\cup B=S,\qquad A\cap B=\varnothing,
\]
且
\[
\forall a\in A,\ \forall b\in B,\ a<b,
\]
则称 $A\mid B$ 是 $S$ 的一个分割。
也就是说，$A\mid B$ 把 $S$ 分成前后两段：
一段是 $A$，一段是 $B$，两段合起来正好是 $S$，且互不重叠。

分割具有以下几条性质：
\begin{enumerate}
\item \textbf{三歧性}：任意的有序数集 $S\subset\mathbb R$，
$\forall a,b\in S$，只有 $a>b$、$a=b$、$a<b$ 三种情况；
\item \textbf{不空不漏}：$A,B\neq\varnothing$，且 $A\cup B=S$；
\item \textbf{不乱}：$\forall a\in A,\ \forall b\in B$，都有 $a<b$；
\item \textbf{不相交}：若 $A\mid B$ 是一组分割，则 $A\cap B=\varnothing$。
\end{enumerate}

\subsubsection*{dedikind 连续性}

若稠密有序集 $S$ 的每个分割 $A\mid B$
都存在唯一的 $c\in S$，使
\[
\forall a\in A,\ \forall b\in B,\quad a\leq c\leq b,
\]
则称 $S$ 具有 dedikind 连续性（即 $S$ 是连续的）。
这里的 $c$ 称为分界点。

上述定义有三个要点：
\begin{enumerate}
\item \textbf{稠密}：$\forall a,b\in S,\ a<b$，
$\exists c\in S,\ a<c<b$；
\item \textbf{存在且唯一}：$\exists!\,c\in S$，即分界点存在且唯一；
\item 分界点 $c\in S$：$A$ 中的数都不超过 $c$，
$B$ 中的数都不小于 $c$，即所有 $a\in A$、$b\in B$ 都被
$c$ 分在两边；而 $c$ 本身必须落在 $S$ 中，
要么在 $A$ 这一边，要么在 $B$ 这一边。
\end{enumerate}

\subsubsection*{$\mathbb Q$ 不具有 dedikind 连续性}

为证明 $\mathbb Q$ 不连续，只需找出 $\mathbb Q$
的一个分割 $A\mid B$，使该分割在 $\mathbb Q$ 内没有分界点。

考虑分割
\[
A=\{x\in\mathbb Q:x<\sqrt2\},\qquad
B=\{x\in\mathbb Q:x\geq\sqrt2\}.
\]
因为 $\sqrt2\notin\mathbb Q$，两集合构成 $\mathbb Q$
的一个分割（$A\cup B=\mathbb Q$ 且 $A\cap B=\varnothing$）。
下面证明 $A$ 中没有最大的数，同理 $B$ 中没有最小的数。

我们的目标是寻找一个 $\epsilon$，使得
\[
\forall a\in A,\ \exists\epsilon>0,\ \epsilon\in\mathbb Q,\quad
a+\epsilon\in A,
\]
即要求
\[
\begin{cases}
\epsilon>0,\\
\epsilon\in\mathbb Q,\\
a+\epsilon<\sqrt2.
\end{cases}
\]

先按放缩法估算。由 $a+\epsilon<\sqrt2$ 得
\[
(a+\epsilon)^2<2
\iff a^2+2a\epsilon+\epsilon^2<2.
\]
若 $\epsilon<2a$，则 $\epsilon^2<2a\epsilon$，于是
\[
a^2+4a\epsilon<2
\iff \epsilon<\frac{2-a^2}{4a}.
\]
故只需取
\[
\epsilon<\min\left\{2a,\ \frac{2-a^2}{4a}\right\},
\]
即可满足要求。
（注意：$\epsilon^2<2a\epsilon$ 与 $\epsilon<2a$ 只有在
$a>0$ 时才等价；而当 $a$ 靠近 $0$ 时，$2a$ 可能已经大于
$\sqrt2-a$，因此还须补上 $\epsilon<\sqrt2-a$ 这一条限制，
方能保证 $a+\epsilon<\sqrt2$，也才保证 $\epsilon>0$。）

下面给出一个更直接的取法。对任意 $a\in A$，
总有 $\sqrt2-a>0$，令
\[
\epsilon=\frac12\min\{1,\ \sqrt2-a\}>0 .
\]
由 $1,\sqrt2\in\mathbb Q$ 及 $a\in\mathbb Q$ 可知
$\epsilon\in\mathbb Q$；并且
\[
\epsilon\leq\frac12(\sqrt2-a)<\sqrt2-a,
\]
从而
\[
a+\epsilon<a+(\sqrt2-a)=\sqrt2,
\]
即 $a+\epsilon\in A$。
这说明 $A$ 中不存在最大的数：
对 $A$ 中任何一个数 $a$，总能在 $A$ 中找到一个比它更大的数。
（同理，对任意 $b\in B$，取
$\epsilon=\frac12\min\{1,\ b-\sqrt2\}$，
可证 $b-\epsilon\in B$，故 $B$ 中不存在最小的数。）

因此，$A$ 中没有最大的数，$B$ 中没有最小的数。
该分割在 $\mathbb Q$ 内没有分界点，
即不存在 $c\in\mathbb Q$ 使
\[
\forall a\in A,\ \forall b\in B,\quad a\leq c\leq b.
\]
所以 $\mathbb Q$ 不具有 dedikind 连续性，$\mathbb Q$ 不连续。

\end{document}
